% =====================================================================
%  Betriebsanweisung Lasercutter LTT iLaser 4000 und Epilog Zing
%  Eigenes PDF: Betriebsanweisung_Lasercutter.tex, Seite in der Einweisung
%  über \BAseite. Layout: fablab-document/fablab-ba.sty
%  Inhalt: Kapitel "Was du dir merken musst" und "Materialien" der Einweisung.
% =====================================================================
\BAsetup{
  typ=maschine,
  titel={Lasercutter LTT iLaser 4000 (\glqq LTT\grqq) und Epilog Zing (\glqq Zing\grqq)},
  taetigkeit={Schneiden und Gravieren, Materialwechsel, Reinigen des Schneidetischs},
  nummer=BA-LC-01,
  stand=TT.MM.JJJJ,                   % Datum der Freigabe
  verantwortlich={Name der verantwortlichen Person},
  bildcode={\includegraphics[width=0.75\BAsymw]{fablab-document/ba-symbole/W004}},
  schrift=\footnotesize,
  entwurf,                            % entfernen, sobald freigegeben
}

\BAkopf

% ---------- 1 Anwendungsbereich --------------------------------------
\BAblock{ANWENDUNGSBEREICH}{M002}{%
  \begin{itemize}
    \item Schneiden und Gravieren \textbf{nur der erlaubten Materialien}
          (Liste in der Einweisung). Was dort nicht steht, ist verboten.
    \item Benutzung nur nach Einweisung; Hinweise des Herstellers beachten.
  \end{itemize}}

% ---------- 2 Gefahren -----------------------------------------------
\BAblock{GEFAHREN FÜR MENSCH UND UMWELT}{W004,W021,W016}{%
  \begin{itemize}
    \item \textbf{Laserstrahlung}: unsichtbarer Infrarot-Laser mit hoher Leistung. Er ist
          nur bei geschlossenem Deckel und geschlossener Frontklappe in Betrieb. Der rote
          Hilfslaser ist harmlos, aber nicht direkt (z.\,B. über einen Spiegel) ins Auge lenken.
    \item \textbf{Brandgefahr}: Acryl, Holz, Papier u.\,ä. können brennen. Brennendes Material
          kann Linse und Lasercutter zerstören.
    \item \textbf{Giftige und ätzende Dämpfe} bei ungeeignetem Material: Salzsäure und Flusssäure
          (z.\,B.\ PVC, Teflon), Blausäure (z.\,B.\ Polyamid, Elastan), Gestank (ABS, Epoxidharz).
    \item \textbf{Verbrennungsgefahr} an heißem Werkstück, Resten und Tischgitter.
    \item \textbf{Quetschgefahr} durch den schnell bewegten Laserkopf und den Arm bei
          geöffnetem Deckel (Testlauf, Zing).
    \item \textbf{Sachschaden} durch Kollision des Laserkopfs oder Magnete am Gerät.
  \end{itemize}}

% ---------- 3 Schutzmaßnahmen ----------------------------------------
\BAblock{SCHUTZMASSNAHMEN UND VERHALTENSREGELN}{P001,P010}{%
  \begin{itemize}
    \item \textbf{Nur erlaubte Materialien}, keine unbekannten Kunststoffe. Keine leicht
          entzündlichen oder gefährlichen Gegenstände (z.\,B. befüllte Feuerzeuge, Handys
          mit Akku).
    \item \textbf{Solange der Laser arbeitet, direkt beim Lasercutter bleiben und zusehen.}
    \item Absaugung und Air Assist müssen laufen (beim LTT automatisch). Nach dem Job kurz
          warten, bis die Dämpfe abgesaugt sind.
    \item Werkstück und Tischgitter nach dem Lasern abkühlen lassen, glimmende Reste nicht wegwerfen.
    \item Nichts verstellen oder aufschrauben, Schutzschalter nicht manipulieren.
          Einzige Ausnahme: die grünen Schrauben des Schneidetischs.
    \item Keine Magnete am oder im Lasercutter. Wabengitter nicht belasten.
    \item Laserkopf muss sich frei bewegen können. Autofokus (LTT) auf den \emph{höchsten}
          Punkt des Werkstücks. Nicht in den Verfahrbereich greifen.
  \end{itemize}}

% ---------- 4 Störungen ----------------------------------------------
\BAblock{VERHALTEN BEI STÖRUNGEN}{W001,F001,P011}{%
  \begin{itemize}
    \item \textbf{Kleine Flamme}: Pause-Taste (LTT) bzw.\ Stop-Taste (Zing) drücken,
          Absaugung und Air Assist prüfen. Der Laser fährt noch einige Linien zu Ende.
    \item \textbf{Große Flamme oder starker Rauch}: \textbf{Deckel leicht öffnen} (Laser schaltet
          ab, der Antrieb fährt weiter: nicht hineingreifen; LTT: \textbf{Not-Aus}, Zing: Kippschalter
          rechts unten). \textbf{Betreuer rufen.}
    \item Löschen nur mit dem \textbf{CO\textsubscript{2}-Löscher} in kurzen Sprühstößen, \textbf{nie mit
          Wasser, Pulver oder Schaum}. Standort: gegenüber vom Laser beim Kassenterminal und am
          Durchgang zur Elektrowerkstatt.
    \item Bei weiterer Eskalation \textbf{Feuerwehr 112}, Personen warnen, Raum verlassen
          (Erstickungsgefahr durch CO\textsubscript{2} und Rauch).
    \item \textbf{Laserkopf stößt an}: sofort Not-Aus, Betreuer informieren. Netzstecker
          ziehen, Hinweiszettel anbringen. \textbf{Nicht zum Test wieder einschalten}, bis die
          Achsen geprüft sind.
  \end{itemize}}

% ---------- 5 Erste Hilfe --------------------------------------------
\BAblock{VERHALTEN BEI UNFÄLLEN – ERSTE HILFE}{E003}{%
  \begin{itemize}
    \item Ruhe bewahren, Gerät abschalten, Betreuer/Ersthelfer verständigen.
    \item Jeden Unfall den Betreuern melden und im Verbandbuch eintragen.
  \end{itemize}\vspace{1.5mm}
  \textbf{Notruf: 112}\qquad FabLab: 09131\,/\,85\,28013}

% ---------- 6 Instandhaltung -----------------------------------------
\BAletzterblock{INSTANDHALTUNG, ENTSORGUNG}{M021}{%
  \begin{itemize}
    \item Wartung, Reinigung der Optik und Filterwechsel nur durch eingewiesene Betreuer,
          bei ausgeschaltetem Gerät (Netzstecker ziehen); Wartungsliste führen.
    \item Vor dem Einschalten nach der Wartung prüfen, dass keine Gegenstände im Laser liegen.
    \item Elektrische Prüfung nach DGUV Vorschrift 3 gemäß festgelegter Prüffrist.
  \end{itemize}}
